% Master thesis proposal: extensions of the 2023 thesis, tested on card data (notes/card-data.md, Step 8).
% Written for a student who knows neither the 2023 thesis nor the repository. Each extension says
% what to derive and points to the handbook (formula and proof idea) to check against.
% Built by CI (xelatex) after the handbook, to /thesis-proposal.pdf. Handbook result numbers come
% from handbook/handbook.aux (package xr), so build the handbook first:
%   (cd handbook && latexmk -xelatex handbook.tex) && (cd studies/thesis-topics && latexmk -xelatex proposal.tex)
\documentclass[10pt]{article}

\usepackage{amssymb, amsmath}
\usepackage[margin=2.1cm]{geometry}
\usepackage{titlesec}
\titlespacing*{\section}{0pt}{8pt}{3pt}
\titlespacing*{\subsection}{0pt}{6pt}{2pt}
\usepackage{booktabs, array}
\newcommand{\rowsep}{\specialrule{0.2pt}{3pt}{3pt}}
\usepackage{xr-hyper}
\usepackage{xcolor}
% Link colours as in the paper and the handbook: internal links in the text colour, external in a muted blue.
\definecolor{linkblue}{RGB}{30,60,110}
\usepackage[colorlinks=true, linkcolor=black, citecolor=black, urlcolor=linkblue]{hyperref}
\externaldocument[H-]{../../handbook/handbook}

\newcommand{\site}{https://mikaelandblom-coder.github.io/credible-lean/}
\newcommand{\repo}{https://github.com/mikaelandblom-coder/credible-lean}
% \hb{label}: the number of a handbook result (no link: the handbook is a separate PDF).
\newcommand{\hb}[1]{\ref*{H-#1}}
\newcommand{\Handbook}{\href{\site handbook.pdf}{Handbook}}
\newcommand{\hmu}{\hat\mu}
\newcommand{\cov}{\operatorname{cov}}
\newcommand{\Var}{\operatorname{Var}}
\newcommand{\bu}{\mathbf 1}

% One extension: \ext{title}{idea}{model}{derive items}{data question}{check against}{size}
\newcommand{\ext}[7]{%
  \subsection*{#1}
  \textit{Idea.} #2\par\smallskip
  \textit{A suggested model.} #3\par\smallskip
  \textit{Derive.}
  \begin{enumerate}\setlength{\itemsep}{1pt}\setlength{\parskip}{0pt}#4\end{enumerate}
  \textit{Data question.} #5\par\smallskip
  {\small\textit{Check against:} Handbook #6 \quad\textperiodcentered\quad \textit{Size:} #7}\par}

% One other approach: \other{title}{idea}{question}{check against}
\newcommand{\other}[4]{%
  \subsection*{#1}
  \textit{Idea.} #2\par\smallskip
  \textit{Question.} #3\par\smallskip
  {\small\textit{Check against:} #4}\par}

\setlength{\parindent}{0pt}
\setlength{\parskip}{4pt}

\title{Correlated credibility: extensions, tested on card data\\[2pt]
  {\large Master thesis proposal in actuarial mathematics}}
\author{\vspace{-2em}}
\date{Draft for discussion, \today}

\begin{document}
\maketitle
\vspace{-1.5em}

\section{Aim}

Credibility theory prices a segment of a portfolio (a \emph{cohort}: a region, a product, an
underwriting year) by blending its own experience with the portfolio's. The classical
Bühlmann--Straub model assumes that cohorts are independent. A master thesis (Andblom, 2023) dropped
that assumption: cohorts may be correlated, as when one storm or one drop in demand hits them all.
Since then the model has been extended in several directions.

This thesis takes some of those extensions, derives their formulas, and tests them on real data.
You derive the formulas yourself, then check them against the handbook (Section~\ref{sec:where}),
which states each formula with a note on its use and what to watch out for. The paper explains the
ideas behind the formulas, and each handbook result links to its machine-checked proof. Deriving is the best way to learn the subject,
and it makes the thesis less dependent on mathematics generated with AI tools. The theory is
general; the data are monthly claims of a travel insurance sold through payment cards
(Section~\ref{sec:cards}), where each extension answers a concrete pricing question.

\textbf{Status.} This is a draft, shared early to discuss the ideas. The formulation itself is
part of the work: review it, and propose changes where the models, the questions or the scope can
be better. The repository behind it is still developing. When the thesis starts, its state will be
frozen as a release, and the thesis will cite that release.

\section{The starting point: credibility for correlated cohorts (2023)}

\textbf{Data.} Cohorts $j = 1, \dots, J$ are observed in periods $t = 1, \dots, T$: a volume
$w_{jt}$ (e.g.\ insured years) and an observed ratio $\mu_{jt}$ (claims per unit of volume). Each
cohort has a true level $\mu_j$, which we want to estimate.

\textbf{Model.} Only first and second moments are assumed, no distribution:
\[
  \mu_{jt} = \mu_j + \sigma_{jt}\,\varepsilon_{jt},
  \qquad \mu_j = \mu + \tau\eta_j,
  \qquad \sigma_{jt}^2 = \frac{\sigma_j^2}{w_{jt}},
  \qquad \cov(\varepsilon_{it}, \varepsilon_{js}) = \rho_{ij}\,\delta_{ts}.
\]
The cohort effects $\eta_j$ and the noise $\varepsilon_{jt}$ are centred with unit variance, and
the cohort effects are uncorrelated with each other and with the noise. The new assumption is the
last one: noise in the same period is correlated across cohorts. With $\rho_{ij} = \delta_{ij}$
(no correlation) and a common $\sigma_j = \sigma$ this is Bühlmann--Straub.

\textbf{Prior art.} This model has a precedent, which a thesis should position itself against.
Poon (2013, an MSc project, read in full) and Poon and Lu (2015) have noise correlated across
risks within a period, the best linear predictor, and the same cross-product estimators; Schnapp
et al.\ (2000) is the original, with equal volumes. The common-effect line (Yeo and Valdez 2006;
Wen, Wu and Zhou 2009; Ebrahimzadeh et al.\ 2012, 2013) instead shares one \emph{permanent} effect
over all periods. A literature check (\href{\repo/blob/main/studies/literature.md}{literature
notes}, \S10q) found these \emph{not} in the earlier work: the credibility factor for any positive
definite $K$ and when it leaves $[0, 1]$, the exact price of the blend form and of ignoring the
correlation, exposure balance, cohorts observed over different periods, and the group model with
correlated noise. A first task is to check that claim against the papers.

\textbf{The recipe.} Every result follows from one matrix.
\begin{enumerate}\setlength{\itemsep}{1pt}
  \item Within a cohort, the volume-weighted mean $\hmu_j = \sum_t w_{jt}\mu_{jt} / w_{j\ast}$ is
    optimal ($w_{j\ast} = \sum_t w_{jt}$).
  \item $K_{ij} = \cov(\hmu_i, \hmu_j) = \sigma_i\sigma_j\rho_{ij}\,S_{ij}/(w_{i\ast}w_{j\ast}) +
    \tau^2\delta_{ij}$ with $S_{ij} = \sum_t \sqrt{w_{it}w_{jt}}$. It must be positive definite.
  \item The portfolio estimate $\hmu = \sum_j b_j\hmu_j$ uses the GLS weights
    $b = K^{-1}\bu / \bu^{\mathsf T}K^{-1}\bu$. A cohort that moves with many others gets less weight.
  \item The blend $z_j\hmu_j + (1 - z_j)\hmu$ has the best credibility factor
    $z_j = (1 - b_j)\tau^2 / (K_{jj} - 1/\bu^{\mathsf T}K^{-1}\bu)$.
  \item $\tau^2$ and $\sigma_i\sigma_j\rho_{ij}$ have unbiased moment estimators, built from
    cross-products of deviations from the cohort means.
\end{enumerate}
Every extension below keeps steps 3--4. It changes \emph{how $K$ is built} (step 2), and sometimes
what is predicted. So each extension asks for the same work: build its $K$, show that it is
positive definite, and find estimators for the new parameters. \Handbook: the model and the recipe
in Section~1, steps 1--5 in Results \hb{res:invVar}, \hb{res:gls}, \hb{res:z},
\hb{res:estimation}.

\section{Which assumptions the extensions relax}

\begin{center}\small
\begin{tabular}{>{\raggedright}p{4.9cm} >{\raggedright}p{5.4cm} >{\raggedright\arraybackslash}p{4.4cm}}
  \toprule
  \textbf{The 2023 assumption} & \textbf{Why it can fail} & \textbf{Extension} \\
  \midrule
  All cohorts observed in all periods & Products start and stop & \textbf{E1 Entry and exit} \\ \rowsep
  Correlation between cohorts falls with volume & A demand shock hits large and small cohorts alike & \textbf{E2 A shock that does not shrink} \\ \rowsep
  The target is a cohort's level, the predictor a blend & One wants next year's total; better linear predictors exist & \textbf{E3 The best linear predictor} \\ \rowsep
  One level per cohort and period & Monthly data have a season & \textbf{E4 Season} \\ \rowsep
  Periods are independent & Demand persists over months and years & \textbf{E5 A persistent common factor} \\ \rowsep
  One ratio per cohort & Frequency and severity behave differently & \textbf{E6 Frequency and severity} \\ \rowsep
  The common movement is unobserved & A public index (e.g.\ passenger numbers) tracks it & \textbf{E7 A known demand index} \\ \rowsep
  Parameters are known & They are estimated from the same, skewed data & \textbf{E8 The price of estimation} \\ \rowsep
  Cohort levels are unrelated & Cards of one country are alike & \textbf{E9 Cards grouped by country} \\
  \bottomrule
\end{tabular}
\end{center}

Each model below is a \emph{suggestion}: change it if the data or the literature call for
something else. The scope is flexible and follows the ambition of the thesis. E3 is a natural
first chapter for E4--E7 and E9; from there one can go deep on one or two extensions or cover more,
and the review of Section~\ref{sec:review} can be a part of its own. Whatever the choice, it needs
the method of Section~\ref{sec:method}. Section~\ref{sec:other} lists further approaches that do not
fit the nine, to consider or to use as comparators.

\section{The extensions}

\ext{E1. Cohorts that enter and leave}
{In 2023 every cohort was observed in every period. In practice products are launched and closed,
and a closed cohort's history still says something about the others.}
{Cohort $j$ is observed in the periods $T_j$, with the 2023 moments on the observed cells.}
{\item $K_{ij}$ when correlation acts only on the periods two cohorts share.
 \item A condition for $K \succ 0$ that involves only cohorts observed at the same time.
 \item An unbiased estimator of $\sigma_i\sigma_j\rho_{ij}$ for a pair that shares only some
   periods: what replaces $T - 2$?}
{Does keeping the history of closed cards improve the prediction for the open ones? The exact
results assume exits unrelated to the claims. What if a card was closed for its results?}
{Results \hb{res:unequal}, \hb{res:unequal_est}}{small}

\ext{E2. A common shock that does not shrink with volume}
{In 2023 the covariance between two cohorts falls like $1/\sqrt{w_{it}w_{jt}}$: double the volume
and the shared shock halves. A real shock (a storm, a drop in travel demand) hits a cohort's ratio
whatever its size.}
{A period shock with covariance $C$ plus process noise,
 $\cov(\mu_{it}, \mu_{js}) = \delta_{ts}\bigl(C_{ij} + \delta_{ij}s_i^2/w_{it}\bigr) + \tau^2\delta_{ij}$.}
{\item $K$ for volume-weighted means, and $K \succ 0$.
 \item Unbiased estimators of $C$ and $s^2$.
 \item Why $s^2$ is identified only when volumes vary over time, and when this model and the 2023
   model coincide.}
{Does the covariance between cards' monthly ratios shrink with volume?}
{Results \hb{res:shock}, \hb{res:shock_est}}{medium}

\ext{E3. The best linear predictor of any target}
{The blend is one linear predictor among many, and it targets the level $\mu_j$. A price is for
next year's total, and the best predictor may weight each cohort separately.}
{Data $x$ with $\cov(x) = K$, a target $Y$ with $\ell = \cov(x, Y)$ and $v = \Var Y$, and unknown
means $Ex = X\beta$, $EY = x_0^{\mathsf T}\beta$.}
{\item For which $d$ is $d^{\mathsf T}x$ unbiased for every $\beta$?
 \item The best such $d^*$ (Lagrange multipliers) and its error. Which part of the error is the
   price of estimating $\beta$?
 \item Any other unbiased $d$ loses exactly $(d - d^*)^{\mathsf T}K(d - d^*)$. When is the 2023
   blend the best predictor?
 \item Next year's outcome instead of the level: why $z$ stays the same.}
{On the data, how much does the best predictor gain over the blend?}
{Results \hb{res:blup}, \hb{res:forecast}}{small to medium}

\ext{E4. Monthly data and the season}
{Monthly data have a season. The usual shortcut is to add up whole years and hope the season drops
out. It does not when volumes are seasonal: a cohort sold mainly in summer leans on summer months.}
{$E\mu_{jt} = \beta_{m(t)}$, one unknown mean per calendar month (or per cohort and month), as fixed
effects in E3; the target is the next 12 months.}
{\item The unbiasedness conditions: one per calendar month.
 \item When the predictor that ignores the season equals the seasonal one, and its bias otherwise.
   Try one cohort, two years of two months, twice the volume in the first month.
 \item With known covariances, aggregating to annual totals never helps. What can change once the
   season must be estimated?}
{Does a monthly fit with an estimated season beat annual totals once everything is estimated?}
{Results \hb{res:blup}, \hb{res:season}}{medium}

\ext{E5. A persistent common factor}
{A period shock is gone the next month. Travel demand persists over months and years, and in the
pandemic it broke. Then past months say something about future ones, through the factor.}
{Add $\lambda_i\lambda_j\Gamma_{ts}$ to E2, with loadings $\lambda_j$ and a covariance $\Gamma$ over
months, e.g.\ AR(1), $\Gamma_{ts} = \gamma\phi^{|t-s|}$, or a block $\Gamma$ with a break at known
dates.}
{\item $K$ and $\ell$ for the next 12 months, and $K \succ 0$ for any observed cells.
 \item A small example where a best weight is negative although $\phi > 0$. Why is it not an error?
 \item Products of two cohorts' 12-month changes at lag $\ge 1$: unbiased for the factor, and free
   of the shock and the season?
 \item Why loadings need three cohorts. When do relative prices between cohorts not depend on the
   factor?}
{Are lagged products between cards different from zero? Does the factor model predict better than
the shock model? Which treatment of the pandemic predicts the later years best?}
{Results \hb{res:factor}, \hb{res:factor_est}, \hb{res:factor_extra}, \hb{res:relativity_factor}}{large}

\ext{E6. Frequency and severity}
{A ratio is frequency times severity, and the two may need different credibility. The levels become
vectors and $z$ becomes a matrix.}
{A level vector (frequency, severity) per cohort, a joint $K$ over cohorts $\times$ components,
frequency weighted by volume and severity by claim counts.}
{\item The best credibility matrix $Z_j$ (a matrix parabola).
 \item Joint and separate prediction coincide when the components do not covary. Is the converse
   true?
 \item The bias of the pure premium as a product of two predictions,
   $E[\hat f\hat s] - E[fs]$.}
{Do frequency and severity residuals covary, and does a joint fit predict better?}
{Results \hb{res:vector}, \hb{res:freqsev_joint}, \hb{res:prodBias}}{medium}

\ext{E7. A known demand index}
{If a public series (passenger statistics) tracks demand, use it as a covariate with a
cohort-specific sensitivity. This is Hachemeister's regression credibility, with demand in place of
time.}
{$\mu_{jt} = \alpha_j + \beta_j D_t + \text{noise}$, with $(\alpha_j, \beta_j)$ random around a
common mean and cohorts correlated as in E2.}
{\item Hachemeister's credibility matrix as an instance of E6's $Z_j$, with correlated cohorts.
 \item Equal sensitivities: when does the index act as a plain exposure adjustment?}
{Given the index, how much correlation between cards remains?}
{Result \hb{res:vector}}{medium}

\ext{E8. The price of estimation}
{Formulas with known $K$ are the oracle. In practice $K$ is estimated from the same data, and claims
are skewed. That can bias the premium and cost accuracy.}
{Any fit whose weights $\hat d$ are computed from the data, against the best weights $d^*$ for the
true moments.}
{\item The error of the plug-in predictor is $\mathrm{mse}(d^*) + EB^2 + 2E[AB]$, with $A$ the
   oracle's error and $B$ the plug-in's distance from it. Under normal noise $E[AB] = 0$ (a result
   that assumes a distribution: say so).
 \item The bias of estimated weights, $\sum_i \cov(\hat d_i, \hmu_i)$: why it is negative for
   right-skewed claims, and why balancing the total premium removes it at portfolio level.
 \item An estimated table need not be positive definite. Projecting it onto valid tables never
   increases the error.
 \item Standard errors of the moment estimators: their variance needs third and fourth moments
   (the general formula is Handbook Result \hb{res:varQuad}).}
{How skewed and how heavy-tailed are the residuals, and what does balancing the total premium cost
in accuracy?}
{Results \hb{res:randomWeights}, \hb{res:plugIn}, \hb{res:valid_tables}}{medium}

\ext{E9. Cards grouped by country}
{Cards come from several countries. Cards of one country share its market, distribution and claims
handling. Card tiers (e.g.\ classic, gold) are products of their country and not comparable across
countries, so cards are nested within countries: hierarchical credibility, with correlated cohorts
on top.}
{A country effect shared by its cards, $\mu_j = \mu + a_{c(j)} + \tau\eta_j$, with uncorrelated
country effects of variance $\sigma_a^2$, so
$\cov(\mu_i, \mu_j) = \sigma_a^2[c(i) = c(j)] + \tau^2\delta_{ij}$; noise as in E2.}
{\item $K$ and $\cov(\hmu_i, \mu_j)$; the best predictor of $\mu_j$ through E3. How much does a
   small card borrow from its own country, and how much from the rest?
 \item The same as two stages: a country estimate, then the portfolio.
 \item Unbiased moment estimators of $\sigma_a^2$ and $\tau^2$, and why the 2023 $\hat\tau^2$ is
   biased when cards share a country level.
 \item Tier within a country as a known difference (a fixed effect) or as part of the card's own
   level? If some tiers do turn out comparable across countries, a crossed tier effect is new
   ground.}
{Do cards of one country resemble each other beyond their own experience?}
{Results \hb{res:phi}--\hb{res:est_groups}}{medium}

\medskip
\textit{Other directions in the handbook, not needed for the card data:} a known claim trend
(Result \hb{res:trend}), when to merge cohorts and how to choose them (Results \hb{res:merge},
\hb{res:pool_intervals}, \hb{res:blend_intervals}), and balance (Results \hb{res:balance}--\hb{res:commonShift}).

\section{Other approaches worth considering}\label{sec:other}

Neither extensions of the model nor part of the 2023 recipe, these are alternatives, baselines and
checks that would strengthen a thesis on the card data. Each is optional.

\other{A1. Covariates first: a model for the means, credibility on the residuals}
{Cards differ in country, tier and channel. Instead of nesting cards in countries (E9), fit a GLM
or another learner on those covariates and use it as the collective mean; credibility then blends
each card's own residual with zero. The weight is the usual factor, from the spread of the
residuals across cards, with no tuning. This is a case of E3, with the learner in the place of
$X\beta$ (``small area learning''). On the log scale it is a credibility correction to a
multiplicative tariff.}
{Does it beat nested levels, and how much does fitting the learner on the same data bias the
premium (cross-fitting is the remedy)? Do tiers behave alike across countries once covariates are
in?}
{Handbook Results \hb{res:blup}, \hb{res:randomWeights}. Reluga et al.\ (2026); Ohlsson (2008);
Nelder and Verrall (1997).}

\other{A2. A likelihood benchmark}
{All alternatives in Section~\ref{sec:method} are moment fits. Add a Poisson or Tweedie mixed
model with correlated random effects, fitted by REML or Laplace (\texttt{lme4}, \texttt{glmmTMB})
or Bayesian. With conjugate priors in an exponential family, linear credibility is the exact Bayes
rule (Jewell 1974), so the question is where it is not: skewed cells, small cells, heavy tails.}
{What does the distribution-free fit give up, if anything, and when? A result that uses a
distribution is fine if it says so.}
{Handbook Result \hb{res:varQuad} (what skewness and kurtosis change); Result \hb{res:randomWeights}.}

\other{A3. Which cards to pool}
{Merge cards that are alike. With known levels the best groups are intervals of the sorted levels
(weighted $k$-means on a line), found by dynamic programming. In practice the levels are estimated
from the data you then fit, so groups look more different than they are: a winner's curse. Cutting
the sorted estimates where an estimated risk is smallest was 2--74\% worse than one group in our
simulations, and with a separate $\tau^2$ per group the optimal groups are not intervals.}
{How large is the curse in general, and is there a correction, as the price of estimation corrects
a fitted model's error? Is the number of groups better chosen by gain against price than by a
likelihood criterion? It is open.}
{Handbook Results \hb{res:merge}, \hb{res:pool_intervals}, \hb{res:blend_intervals},
\hb{res:bsPrice}; Findings 14, 41.}

\other{A4. Baselines for the persistent factor}
{A practitioner weights a card's own years by $\lambda^{\text{age}}$ (exponential smoothing; Morris
2026, with $\lambda$ depending on account size). The factor of E5 has a state-space form, so the
best predictor can be computed by a Kalman filter on long histories (Gerber and Jones 1975;
Zehnwirth 1985), and Pinquet (2020) restricts the weights to be nonnegative. In the repository's
simulated monthly world, recency weighting lost 6--31\% against the best predictor when the factor
persists, because the best predictor reads the factor from the \emph{other} cards.}
{Is that the case on the data? Does the best $\lambda$ fall with a card's volume, as the model
predicts? Is a restriction to nonnegative weights worth its cost?}
{Handbook Result \hb{res:blup} (the exact loss of any aggregate predictor), Results \hb{res:factor}, \hb{res:factor_est}.}

\other{A5. The pandemic as an outlier}
{E5 handles it as a break in the factor's covariance at known dates. Alternatives: drop the months,
or use robust credibility, which limits the pull of extreme observations (Künsch 1992; winsorized
losses with quadratic terms, Zhao and Poudyal 2026). Robust versions need moments beyond the
second, so they leave the distribution-free core.}
{Which treatment predicts the later years best, and does it change the cards' ranking?}
{Handbook Results \hb{res:factor}, \hb{res:varQuad}.}

\other{A6. Uncertainty per card}
{The mean squared errors in the handbook average over the cards' levels. For one card whose level
lies far from the portfolio, the shrunk estimate is biased towards the portfolio, and an interval
built from the average error can cover far less than it claims (2--4\% in a small-area study).
The error given the card's level, and its unbiased estimate, are known.}
{Back-test interval coverage per card, not only squared error: does the average-error interval
under-cover the outlying cards, and does the given-level one repair it?}
{Handbook Result \hb{res:givenLevel}; Result \hb{res:plugIn} (quote a fitted error as typical).}

\section{Results to review}\label{sec:review}

The continuation found results that seem to be new, often gaps in the 2023 thesis itself. For any
of them, check: is it correct, is it already known (a reference is as good an outcome as a
confirmation), and does it matter in practice? Each is in the \Handbook{}, whose appendix \emph{Findings} lists all 41 findings with their
counterexamples and Lean proofs, and in \href{\repo/blob/main/FINDINGS.md}{FINDINGS.md} with the
details.
\begin{enumerate}\setlength{\itemsep}{1pt}
  \item \textbf{The 2023 model makes a common shock shrink with volume} (Finding 23, Result
    \hb{res:shock}). With a fixed $\rho_{ij}$ the shared shock falls like $1/\sqrt{w_{it}w_{jt}}$;
    the fix is E2.
  \item \textbf{Whole years do not remove the season} when volumes are seasonal (Finding 24,
    Result \hb{res:season}): one cohort, two years of two months and twice the volume in the
    first month give weights $\tfrac43$, $\tfrac23$ and a bias $(\beta_1 - \beta_2)/3$.
  \item \textbf{Positive autocorrelation does not keep the best weights positive} (Finding 25,
    Result \hb{res:factor}): with an AR(1) factor and $\phi = \tfrac12$ a weight of $-\tfrac1{30}$
    is optimal.
  \item \textbf{Estimated weights under-price skewed portfolios} (Finding 35, Result
    \hb{res:randomWeights}): the bias is exactly $\sum_i\cov(\hat d_i, \hmu_i)$, negative when a
    high draw also raises the cohort's estimated variance; balancing on exposure removes it for
    the portfolio. Compare with inverse-variance weighting in meta-analysis.
  \item \textbf{Optimal groups of cohorts are intervals} (Finding 14, Results
    \hb{res:pool_intervals}, \hb{res:blend_intervals}): pooling cohorts is weighted $k$-means in
    one dimension, and with a common variance the same holds for the blend.
  \item \textbf{The own-data weight stays bounded} where $z_j$ explodes (Finding 17, Result
    \hb{res:z}): near the two-cohort singularity $z_j$ is large or negative, but
    $c_j = b_j + z_j(1 - b_j)$ is not.
  \item \textbf{Which of these are new?} The prior art of Section~2 contains the model and its
    estimators. Check each claim of novelty above against Poon and Lu (2015) and Schnapp et al.\
    (2000) before building on it, and write down what each paper has and lacks.
\end{enumerate}

\section{Comparing models on data}\label{sec:method}

The theory says what is best \emph{if} a model holds, not which model holds. Decide on the data by
out-of-sample prediction. Use a rolling origin with no look-ahead: at each month, refit every
alternative on the past only, predict each cohort's next 12 months, and score the squared error.
Compare two alternatives by their \emph{paired} loss differences, so the target year's common
shock cancels, with a Diebold--Mariano test, a HAC variance with lag 11 and the
Harvey--Leybourne--Newbold correction. Do not use shuffled $K$-fold: the common factor correlates
held-out months with the training months.

First check the power in simulation: in a world the size of the data, how often does the test find
a real difference? (The app's first answer: with 6--8 years a 30\% gap is found in only 15--23\% of
draws.) ``No evidence of a difference'' is then a result in itself, and so is finding that a
statement in the handbook is wrong or of little practical use. Score squared error and also interval
coverage per card (A6). With so little power, \emph{combining} fits may beat choosing one
(forecast combination): two predictors' best blend is the credibility factor itself, with
correlated errors. This is an idea, not tested.

\section{The application: travel insurance on payment cards}\label{sec:cards}

A travel insurance is sold through several card products. For each card product (a cohort) and
month there are cards in force, claim counts and paid amounts. The task is to price each card for
the next underwriting year. Classical credibility treats the cards as independent, but travel
demand moves them all at once, with the pandemic as the extreme case. The business is seasonal,
volumes differ widely between cards, cards start and stop within the window, and they come from
several countries, each with its own tiers. The pricing questions map to the extensions:
\begin{itemize}\setlength{\itemsep}{1pt}
  \item Does the correlation between cards shrink with volume, or is it a common shock? (E2, E5)
  \item Does the common movement persist, do cards load on it equally, and does a public index
    capture it? (E5, E7)
  \item Monthly data with an estimated season, or annual totals? (E3, E4)
  \item Frequency and severity combined, separate or joint? (E6)
  \item How should new and closed cards enter the fit? (E1)
  \item How much does a card borrow from the other cards of its country? (E9)
  \item Do country and tier work better as covariates than as nested levels? Which cards can be
    pooled? (A1, A3)
\end{itemize}
\textit{Left to the student:} the number of factors, the treatment of the pandemic, reporting delay
(simplest: use developed months only), and whether new cardholders behave differently.
\textit{Data and publication:} aggregated or anonymized company data by agreement; the thesis text
is public. The analysis must also run on simulated data (the app can produce it), so it can be
reproduced without the company data.

\section{How to work}

\begin{itemize}\setlength{\itemsep}{1pt}
  \item \textbf{Derive first, then check.} Derive the listed results by hand, then compare with the
    handbook, which gives the formula and a practical note; the paper explains why it holds, and
    the proof is behind the handbook's Lean link. Where you disagree, find out who is right.
  \item \textbf{Guidance, not authority.} The continuation was developed with extensive AI
    assistance. Its proofs are machine-checked, so a stated result follows from its definitions. A
    machine cannot check whether the definitions say what is intended, whether a model fits real
    data, or whether a result is already known.
  \item \textbf{No Lean needed.} Every handbook result links to its machine-checked proof in the
    proof assistant Lean. Reading it is optional.
  \item \textbf{Your own code} (R or Python) for estimation and back-testing. The app is for
    intuition and for checking the method, not a dependency.
  \item \textbf{Moments, not distributions.} The core model assumes first and second moments only.
    A result that assumes a distribution (normal noise, gamma amounts) is fine if stated as such.
  \item \textbf{No hand-set tuning.} Every parameter of a fitted model is estimated from the data.
\end{itemize}

\section{Where to find things}\label{sec:where}

\begin{center}\small
\begin{tabular}{>{\raggedright}p{3.6cm} >{\raggedright\arraybackslash}p{11.2cm}}
  \toprule
  \href{\site handbook.pdf}{Handbook} & Every formula, with how to use it and what to watch out for, linked to its Lean proof;
    the appendix \emph{Findings} lists all 41 findings with counterexamples. Result numbers in this
    proposal refer to it. \\ \rowsep
  \href{\site paper/}{Paper} & The up-to-date version of the 2023 thesis: concepts, intuition and the reasons behind the
    formulas, organized for lookup, with a literature map. \\ \rowsep
  \href{\site app/}{App} & Simulated worlds: draw data, fit models that never see the truth,
    compare with the best possible fit. ``Monthly cohorts'' is a world like the card data. \\ \rowsep
  \href{\repo}{Repository} & Everything below, with its history. \\ \rowsep
  \href{\repo/blob/main/FINDINGS.md}{FINDINGS.md} & Gaps and surprises, each with a counterexample
    or proof. \\ \rowsep
  \href{\repo/blob/main/studies/literature.md}{Literature notes} & Credibility, small-area
    estimation, covariance shrinkage; \S10 for monthly card data. \\ \rowsep
  \href{\site}{Blueprint} (optional) & Every statement with its full proof and a dependency graph. \\
  \bottomrule
\end{tabular}
\end{center}

\section{Suggested phases}

\begin{enumerate}\setlength{\itemsep}{1pt}
  \item \textbf{Background.} Read the 2023 thesis (or the paper) and Bühlmann--Gisler ch.~4;
    rederive Section~2 above. Review this proposal.
  \item \textbf{Choose} extensions and derive their results, checking against the handbook.
  \item \textbf{Data.} Descriptive analysis: volumes, season, correlations between cards by lag,
    differences between countries, the pandemic.
  \item \textbf{Simulation check.} Does the comparison of Section~\ref{sec:method} have the power to
    separate your alternatives?
  \item \textbf{Comparison on data} and recommendations for pricing card products, with a list of
    statements checked, corrected or found already known.
\end{enumerate}

\textbf{Prerequisites.} Linear and mixed models (variance components), credibility theory, time
series basics, simulation, R or Python.

\section*{Starting references}
{\small
\begin{itemize}\setlength{\itemsep}{0pt}
  \item Andblom, M. (2023). \emph{Generalized Bühlmann--Straub credibility theory for correlated
    data}. Master Thesis in Actuarial Mathematics 2023:2, Stockholm University.
  \item \emph{credible-lean: correlated credibility theory, formalized in Lean} (2026),
    \url{\repo}. A release will be frozen when the thesis starts.
  \item Bühlmann, H. and Gisler, A. (2005). \emph{A Course in Credibility Theory and its
    Applications}. Springer (chapters 4, 6, 8 and 9).
  \item Goldberger, A. S. (1962). Best linear unbiased prediction in the generalized linear
    regression model. \emph{JASA} 57, 369--375.
  \item Hachemeister, C. A. (1975). Credibility for regression models with application to trend. In
    \emph{Credibility: Theory and Applications}, Academic Press.
  \item Frees, E. W. (2004). \emph{Longitudinal and Panel Data}. Cambridge University Press.
  \item Rao, J. N. K. and Molina, I. (2015). \emph{Small Area Estimation}, 2nd ed. Wiley.
  \item Diebold, F. X. and Mariano, R. S. (1995). Comparing predictive accuracy. \emph{JBES} 13,
    253--263.
  \item Poon, J. (2013). \emph{Spatial cross-sectional credibility models with general dependence
    structure among risks}. MSc project, Simon Fraser University. Poon, J. and Lu, Y. (2015). A
    spatial cross-sectional credibility model with dependence among risks. \emph{North American
    Actuarial Journal} 19(4). Schnapp et al.\ (2000), \emph{CAS Forum}.
  \item Reluga, K., Salvati, N., Ranjbar, S., Tzavidis, N. and van der Laan, M. (2026). Small area
    learning. arXiv:2610.15953. Ohlsson, E. (2008). Combining generalized linear models and
    credibility models in practice. \emph{Scand.\ Actuar.\ J.} 2008(4). Nelder, J. A. and Verrall, R. J.
    (1997). Credibility theory and generalized linear models. \emph{ASTIN Bull.} 27(1).
  \item Jewell, W. S. (1974). Credible means are exact Bayesian for exponential families.
    \emph{ASTIN Bull.} 8(1), 77--90.
  \item Gerber, H. U. and Jones, D. A. (1975). Credibility formulas of the updating type.
    \emph{Trans.\ Soc.\ Actuaries} 27. Zehnwirth, B. (1985). Linear filtering and recursive
    credibility estimation. \emph{ASTIN Bull.} 15(1), 19--36. Pinquet, J. (2020). Poisson models
    with dynamic random effects and nonnegative credibilities per period. \emph{ASTIN Bull.} 50(2),
    585--618. Morris, J. (2026). Logistic credibility with temporal decay. arXiv:2606.08692.
  \item Künsch, H. R. (1992). Robust methods for credibility. \emph{ASTIN Bull.} 22(1), 33--49.
    Zhao, Q. and Poudyal, C. (2026). Robust quadratic credibility. \emph{Ann.\ Actuar.\ Sci.} 20(2),
    302--332.
\end{itemize}}

\end{document}
